\section{Mechanism Layer and Correspondence}\label{supp:mechanisms}
The common inclusion criterion is an original result about an execution-dependent signal, necessary observation access, a detection/transfer condition or a protection action that can be mapped to an Agent setting. The sixteen communication anchors cover interactive Web state leakage, distributed request tracing, IoT/voice event detection, deployment assumptions, residual features and transport cover. Six serving anchors are included by the same criterion: Time Will Tell, Whisper Leak, Early Bird, CachePrune, Weiss et al. and Carlini--Nasr. The latter group includes inspected preprints. All twenty-two have separate locators and transfer limits in \texttt{mechanism-sources.csv}. The main transfer matrix groups the selected relations by the condition needed for an Agent application; it does not turn the anchor set into additional Agent replications.

The original communication extension used twelve recorded discovery/bibliographic queries and selected twelve anchors from fourteen entered leads. WAP5 and vPath require host observations; their graph reconstruction is not a passive encrypted-link result. A retrieved serverless data-flow lead was excluded because it tracks sensitive data with a different observation model. The search does not prove the absence of a passive microservice or serverless attack. Serving mechanisms were integrated from the already inspected context sources, not counted as a second independent discovery sample.

\begin{table*}[t]\centering\small
\caption{Mechanism locators and the observation relationship used in the SoK.}
\label{tab:mature-locators}
\begin{tabularx}{\textwidth}{@{}P{.27\textwidth}P{.22\textwidth}Y@{}}\toprule
Source & Locator & Correspondence boundary \\\midrule
Chen et al.~\cite{chen2010web} & III.B--C & Successive Web-state observations constrain inputs. \\
Chapman--Evans~\cite{chapman2011blackbox} & Sec. 3 & Browser exploration supplies labels and actions. \\
WAP5~\cite{reynolds2006wap5} & Secs. 1--3 & Participating-host socket traces support path inference. \\
vPath~\cite{tak2009vpath} & Secs. 2.2, 3 & Thread/network observations and request--reply assumptions. \\
PINGPONG~\cite{trimananda2020pingpong} & IV.B, V.B--C & Stream signature detection and negative traffic. \\
Ahmed et al.~\cite{ahmed2023voice} & Secs. 3.1--3.2, 4, 7 & Invocation localization before command recognition. \\
Juarez et al.~\cite{juarez2014critical} & Secs. 3--6 & Environment, client and open-world assumptions. \\
Peek-a-Boo~\cite{dyer2012peekaboo} & Sec. I; Fig. 1 & Coarse time/volume features remain useful under padding. \\
Mathews et al.~\cite{mathews2023critical} & Sec. I, findings & Adapted attacks and matched costs. \\
STP~\cite{apthorpe2019stp} & Secs. 6.1, 7.1 & Activity/cover shaping; tunnel endpoints drop cover. \\
Walkie-Talkie; FRONT~\cite{wang2017walkietalkie,gong2020front} & Secs. 3--4; Sec. 4 & Page-load shaping mechanisms need an Agent workload correspondence. \\
Time Will Tell~\cite{timewilltell2024} & Secs. III, VI & Exact token counts and latency proxies are different views. \\
Whisper Leak~\cite{whisperleak2025} & Secs. 3.1, 6.1--6.3 & TLS size/gap inference; new captures and transformed traces differ. \\
Early Bird~\cite{earlybird2025} & II-C, III, V-A--V-B & Prefix and semantic-cache defenses have distinct costs. \\
CachePrune~\cite{wu2026cacheprune} & Secs. 3.1, 4, 5.7 & Filter-disabled component cost does not price enabled protection. \\
Weiss et al.~\cite{weiss2024prompt} & Sec. 2; attack construction & Encrypted streaming token-length sequences. \\
Carlini--Nasr~\cite{carlini2024timing} & Secs. II--III & Optimization-dependent timing; passive/active access separated. \\\bottomrule
\end{tabularx}\end{table*}

Mechanism anchors never enter Agent-source support counts. They establish why the missing capability test is relevant, not its result for an Agent. Transport specifications likewise state what MCP, HTTP/2 and TLS support without demonstrating the privacy of a particular deployment. The artifact records the seven correspondence groups and their setting-level links.

The transfer matrix uses four relationship labels. Retained applies to a classical obligation that remains relevant, such as accounting for unknown boundaries or protected-condition attack adaptation. Extended names an extra execution condition, such as following a tool result into a later model request or tracking the original-task effects of executed cover. Access mismatch rejects a proposed correspondence between host instrumentation and passive packet observation. Untested identifies a performance claim for which the mapped Agent evidence supplies no matched test. A row's relation and the number/publication status of its sources are separate attributes. In particular, STP's tunnel terminator remains valid for Agent deployments that actually use it; WRA's executed research is a different intervention. Classical work also includes multiple observation locations and cross-device dependencies, so a universal single-link assumption is not part of this synthesis.


The focused WF extension adds four published anchors. Wang--Goldberg~\cite{wang2016realistic}, Sections 3--5, supplies the realistic splitting/background precedent. k-FP~\cite{hayes2016kfp}, Sections 4--5, supplies the open-world/feature reference. BAPM~\cite{guan2021bapm} is checked through its official ACSAC abstract and author slides; no quantitative result is imported. ARES~\cite{deng2023ares} is identified by its IEEE publication record and author repository; the inspected extended threat model is arXiv:2501.12622v1, Sections III--IV. That extension is not silently substituted for the 2023 experimental results. These additions refine the localization/transfer correspondence without changing Agent-source support counts.
